% ─────────────────────────────────────────────────────────────────
% CENG 467 Midterm – Question 2: Named Entity Recognition
% Student ID: 310201050
% ─────────────────────────────────────────────────────────────────

\section{Question 2: Named Entity Recognition}

\subsection{Dataset and Annotation Structure}

The CoNLL-2003 English shared-task corpus \cite{sang2003conll} is used. It
consists of newswire articles from the Reuters news collection annotated
with four entity types: Persons (\textsc{PER}), Organisations
(\textsc{ORG}), Locations (\textsc{LOC}) and Miscellaneous entities
(\textsc{MISC}). Annotations follow the BIO scheme (\textsc{B}-prefix
for the first token of an entity, \textsc{I}-prefix for inside tokens, and
\textsc{O} for non-entity tokens), yielding a 9-class tag set:
$\{$O, B-PER, I-PER, B-ORG, I-ORG, B-LOC, I-LOC, B-MISC, I-MISC$\}$.

The corpus is loaded from a static mirror because the HuggingFace
\texttt{datasets} library (v3.0+) no longer supports script-based
loaders and all official mirrors of CoNLL-2003 are script-based. The
downloaded files are normalised to BIO tagging by an idempotent
IOB1$\rightarrow$BIO conversion routine. The standard splits are then
used as released: 14{,}041 training sentences, 3{,}250 validation
sentences, and 3{,}453 test sentences. The test set is used
\emph{exactly once} for final evaluation.

\subsection{Token--Label Alignment}

For the BiLSTM-CRF the alignment is trivial because the corpus is already
pre-tokenised at the word level. For DistilBERT, however, each
pre-tokenised word is further decomposed into one or more WordPiece
sub-tokens. We adopt the canonical alignment scheme: \emph{the first
sub-token of each word inherits the original BIO label, while subsequent
sub-tokens are assigned the special ignore index $-100$}. The
\textsc{CrossEntropyLoss} skips these positions, and during evaluation
only the first sub-token of each word contributes to the predicted
sequence, ensuring that entity-level metrics are computed against
word-aligned predictions.

\subsection{Models}

\subsubsection{BiLSTM-CRF (Hybrid)}

A single-layer bidirectional LSTM with a Conditional Random Field
decoding layer \cite{lample2016neural} is trained from scratch. The
CRF layer explicitly models tag-transition probabilities, which
discourages illegal sequences such as \texttt{O B-PER I-LOC}.

\begin{table}[h]
\centering
\caption{BiLSTM-CRF Hyperparameters}
\label{tab:bilstm_crf_hp}
\begin{tabular}{ll}
\toprule
Hyperparameter & Value \\
\midrule
Vocabulary size           & $10{,}951$ words (min freq.\ 2) \\
Embedding dimension       & 96 \\
Hidden dimension          & 192 \\
Number of LSTM layers     & 1 (bidirectional) \\
Dropout                   & 0.40 \\
Batch size                & 32 \\
Epochs                    & 6 \\
Optimiser / LR            & Adam / $8 \times 10^{-4}$ \\
Weight decay              & $10^{-5}$ \\
Gradient clipping         & 1.0 \\
\bottomrule
\end{tabular}
\end{table}

\subsubsection{DistilBERT (Contextual)}

\texttt{distilbert-base-cased} is fine-tuned for token classification.
The \emph{cased} variant is preferred because capitalisation is a strong
NER feature in English news text (\textit{Apple} vs.\ \textit{apple}).

\begin{table}[h]
\centering
\caption{DistilBERT Fine-Tuning Hyperparameters}
\label{tab:bert_ner_hp}
\begin{tabular}{ll}
\toprule
Hyperparameter & Value \\
\midrule
Pretrained model        & \texttt{distilbert-base-cased} \\
Maximum sequence length & 128 \\
Batch size              & 16 \\
Epochs                  & 3 \\
Learning rate           & $3 \times 10^{-5}$ \\
Weight decay            & 0.01 \\
Warmup ratio            & 10\% \\
Gradient clipping       & 1.0 \\
Scheduler               & Linear with warmup \\
\bottomrule
\end{tabular}
\end{table}

\subsection{Evaluation Results}

Entity-level Precision, Recall and F1 are computed using the
\texttt{seqeval} library, which evaluates \emph{exact span matches}: a
prediction is correct only if both the boundaries and the entity type
match the gold annotation (Table~\ref{tab:q2_results}).

\begin{table}[h]
\centering
\caption{CoNLL-2003 Test-Set Results (entity-level, seed 3102)}
\label{tab:q2_results}
\begin{tabular}{lccc}
\toprule
Model & Precision & Recall & F1 \\
\midrule
BiLSTM-CRF         & 0.8010 & 0.6234 & 0.7011 \\
DistilBERT (cased) & \textbf{0.8895} & \textbf{0.9056} & \textbf{0.8975} \\
\bottomrule
\end{tabular}
\end{table}

The contextual DistilBERT outperforms the from-scratch BiLSTM-CRF by
nearly 20 F1 points, the largest margin observed across all five
questions in this study. A per-class breakdown is given in
Table~\ref{tab:q2_per_class}.

\begin{table}[h]
\centering
\caption{Per-Entity F1 Scores on the Test Set}
\label{tab:q2_per_class}
\begin{tabular}{lcc}
\toprule
Entity Class & BiLSTM-CRF F1 & DistilBERT F1 \\
\midrule
\textsc{PER}  & 0.7258 & \textbf{0.9507} \\
\textsc{LOC}  & 0.7777 & \textbf{0.9245} \\
\textsc{ORG}  & 0.6258 & \textbf{0.8754} \\
\textsc{MISC} & 0.6311 & \textbf{0.7682} \\
\bottomrule
\end{tabular}
\end{table}

The relative gap between models is largest for \textsc{ORG} and
\textsc{PER}, which contain the highest density of rare proper nouns,
and smallest for \textsc{MISC}, which is intrinsically the noisiest
entity class.

\subsection{Error Analysis}

The 5{,}644 gold entities in the test set produced the following
DistilBERT error breakdown:

\begin{table}[h]
\centering
\caption{DistilBERT Error Categories on the Test Set}
\label{tab:q2_errors}
\begin{tabular}{lr}
\toprule
Category & Count \\
\midrule
True Positives  & 5{,}106 \\
Type confusion  & 250 \\
False positives & 236 \\
False negatives & 217 \\
Boundary errors & 71 \\
\bottomrule
\end{tabular}
\end{table}

Three observations follow from this distribution.

\paragraph{Boundary errors are remarkably rare.}
Only 71 of the 5{,}644 gold spans (1.3\%) suffered from boundary
mismatches, all but one of which involved \textsc{MISC} entities (e.g.\
predicting only the head of a multi-word entity). The CRF layer of the
BiLSTM-CRF model was originally introduced to suppress exactly this
class of error, but DistilBERT's softmax classifier achieves a
comparable result through context alone.

\paragraph{Type confusion is the dominant error mode.}
With 250 cases (4.4\% of the gold spans), type confusion is by far the
most common DistilBERT error. Inspection of the recorded examples
reveals classic CoNLL ambiguities: country names used metonymically for
national teams (\textsc{LOC} tagged as \textsc{ORG}), surnames identical
to place names (\textsc{LOC} vs.\ \textsc{PER}), and ambiguous
sentence-initial proper nouns. These errors require world knowledge
that cannot be extracted from the 14k-sentence training set alone.

\paragraph{False positives concentrate in \textsc{MISC}.}
All three of the recorded false-positive examples are \textsc{MISC}
predictions over tokens that are not entities at all. \textsc{MISC} is a
catch-all class for nationality adjectives, event names, and other
``odd'' proper nouns; the model has learnt to over-predict it whenever
capitalisation cues are present without strong contextual constraints.

\subsection{Discussion: Role of Contextual Embeddings}

The BiLSTM-CRF must learn word representations from scratch using only
the 14k-sentence training corpus. Out-of-vocabulary words are mapped to
a single \texttt{<UNK>} embedding, which severely limits generalisation
to rare or unseen proper nouns --- precisely the population that NER
must handle well. Even with a structured CRF decoder enforcing label
consistency, the model achieves only $0.7011$ F1.

DistilBERT's contextual sub-word embeddings address both limitations
simultaneously. WordPiece tokenisation eliminates the closed-vocabulary
problem entirely: any character sequence is decomposable into known
sub-units, allowing the model to construct meaningful representations
for \textit{Schmidt}, \textit{Kyiv}, or \textit{Chrysler} without ever
having seen those tokens during pre-training. Self-attention provides
global context for every position, so the same surface form can receive
different representations depending on its context (e.g.\
\textit{Washington} as a \textsc{PER} versus a \textsc{LOC}). These two
properties combined explain the almost 20-point F1 gap between the two
models.

The principal trade-off is computational: DistilBERT requires roughly an
order of magnitude more memory and compute than the BiLSTM-CRF for both
training and inference, even in its distilled form. Where the latency
or memory budget is tight, BiLSTM-CRF remains a reasonable default;
where accuracy is paramount, contextual representations are clearly
superior.